%==============================================================================
% Makalah Gemastik — Buletin Pagelaran Mahasiswa Nasional Bidang TIK
% Judul: [ISI JUDUL DI SINI]
% Kompilasi: pdflatex nama_file.tex   (jalankan 2x untuk referensi/label)
%==============================================================================
\documentclass[10pt,a4paper,twocolumn]{article}

\usepackage[T1]{fontenc}
\usepackage[utf8]{inputenc}
% Catatan: paket babel dengan opsi [indonesian] TIDAK dipakai di sini karena
% banyak instalasi TeX Live (termasuk lingkungan Claude) tidak menyertakan
% berkas bahasa indonesian.ldf secara default, yang membuat kompilasi gagal
% total. Babel bahasa Indonesia hanya memengaruhi penomoran/nama otomatis
% (mis. label caption bawaan), yang sudah kita override manual di bawah,
% jadi tidak dibutuhkan secara fungsional.
\usepackage{mathptmx}            % Times New Roman untuk teks & matematika
\usepackage[scaled=0.92]{helvet} % Helvetica Type1 asli untuk judul/heading
                                  % (JANGAN pakai \sffamily tanpa helvet: itu
                                  % jatuh ke font bitmap Computer Modern Sans
                                  % yang butuh generate metafont on-the-fly
                                  % dan bisa GAGAL total di lingkungan tanpa
                                  % akses tulis/metafont, termasuk di sini)
\usepackage{amsmath,amssymb}
\usepackage{graphicx}
\usepackage{booktabs}
\usepackage{array}
\usepackage{ragged2e}
\usepackage{caption}
\usepackage{titlesec}
\usepackage{fancyhdr}
\usepackage{url}
\usepackage[hidelinks]{hyperref}
\usepackage{tikz}                % dipakai untuk kotak placeholder gambar
\usepackage[a4paper,left=1.8cm,right=1.8cm,top=2.6cm,bottom=2.2cm,%
            columnsep=0.7cm,headsep=0.5cm]{geometry}

%------------------------------------------------------------------------------
% Penomoran bagian: I., A. (gaya IEEE)
%------------------------------------------------------------------------------
\renewcommand{\thesection}{\Roman{section}}
\renewcommand{\thesubsection}{\Alph{subsection}}

\titleformat{\section}
  {\normalfont\small\bfseries}{\thesection.}{0.5em}{\MakeUppercase}
\titlespacing*{\section}{0pt}{1.4ex plus .3ex}{0.8ex}

\titleformat{\subsection}
  {\normalfont\small\bfseries\itshape}{\thesubsection.}{0.5em}{}
\titlespacing*{\subsection}{0pt}{1.2ex plus .2ex}{0.6ex}

% Heading level 3 (subsubsection): italic, angka Arab + ")", konten paragraf baru
\titleformat{\subsubsection}
  {\normalfont\small\itshape}{\arabic{subsubsection})}{0.5em}{}
\titlespacing*{\subsubsection}{0pt}{1.0ex plus .2ex}{0.4ex}

%------------------------------------------------------------------------------
% Gaya caption
%   - Gambar: label+judul DI BAWAH gambar (default LaTeX figure), Helvetica-ish
%   - Tabel: label+judul DI ATAS tabel (default LaTeX table), Small Caps 8pt
%------------------------------------------------------------------------------
\captionsetup{font=small,labelsep=period,justification=centering}
% Catatan: aturan resmi meminta label+judul tabel dalam "Small Caps".
% Kombinasi shape sc + ukuran footnotesize memicu pembuatan font bitmap
% on-the-fly (mktexpk) yang bisa GAGAL di lingkungan tanpa akses
% tulis/metafont (termasuk lingkungan Claude ini). Untuk keandalan
% kompilasi di semua lingkungan, judul tabel memakai footnotesize+bold
% biasa alih-alih true small caps.
\captionsetup[table]{font=footnotesize,labelfont=bf,skip=3pt}
\captionsetup[figure]{font=small,skip=4pt}

%------------------------------------------------------------------------------
% Header banner + footer (wajib tampil di tiap halaman)
%------------------------------------------------------------------------------
\setlength{\headheight}{26pt}
\pagestyle{fancy}
\fancyhf{}
\renewcommand{\headrulewidth}{0.8pt}
\renewcommand{\footrulewidth}{0pt}
\fancyhead[L]{%
  \footnotesize\bfseries BULETIN PAGELARAN MAHASISWA NASIONAL BIDANG TEKNOLOGI INFORMASI DAN KOMUNIKASI\\%
  \mdseries\scriptsize ISSN: xxxxxx}
% GANTI baris footer di bawah: empat kata pertama judul + nama ketua tim,
% dan volume/nomor/bulan/tahun sesuai edisi yang berlaku.
\fancyfoot[L]{\footnotesize Ketua Tim: Empat kata pertama judul makalah diikuti titik-titik (\ldots)}
\fancyfoot[R]{\footnotesize Volume 1 Nomor 1 [Bulan] [Tahun]}

\setlength{\parindent}{1.2em}
\setlength{\parskip}{0pt}

%------------------------------------------------------------------------------
% Helper: kotak placeholder untuk gambar yang belum tersedia.
% Pakai ini alih-alih \includegraphics kalau file gambar asli belum ada,
% supaya PDF tetap bisa dikompilasi dan placeholder jelas terlihat.
%
% Pemakaian:
%   \placeholderfig{width}{height}{Keterangan singkat isi gambar}
% Contoh:
%   \placeholderfig{0.9\linewidth}{4cm}{Diagram alur sistem: input -> model -> output}
%------------------------------------------------------------------------------
\newcommand{\placeholderfig}[3]{%
  \begin{tikzpicture}
    \draw[dashed, thick, gray] (0,0) rectangle (#1,#2);
    \node[align=center, text width=0.9\dimexpr#1, font=\footnotesize\itshape, gray]
      at (0.5*#1,0.5*#2) {[PLACEHOLDER GAMBAR]\\#3};
  \end{tikzpicture}%
}

\begin{document}

%==============================================================================
% JUDUL, PENULIS, INTISARI (lebar penuh, di atas kolom ganda)
%==============================================================================
\twocolumn[{%
\begin{@twocolumnfalse}
\thispagestyle{fancy}
\vspace{0.5em}
\begin{center}
  % JUDUL: maksimum 12 kata, Helvetica 20pt bold (mathptmx tidak punya
  % Helvetica asli, \sffamily dipakai sebagai pendekatan visual terdekat)
  {\LARGE\sffamily\bfseries [ISI JUDUL MAKALAH — MAKSIMUM 12 KATA]\par}
  \vspace{1.2em}
  % PENULIS: nama keluarga di akhir tiap nama; TANPA gelar/jabatan
  {\normalsize\sffamily\bfseries Nama Ketua Tim\textsuperscript{1}, Nama Anggota 1\textsuperscript{2},
   Nama Anggota 2\textsuperscript{3}, Nama Pembimbing\textsuperscript{4}\par}
  \vspace{0.6em}
  {\footnotesize
   \textsuperscript{1}Nama Fakultas, Nama Institusi, Kota, Kodepos, Negara,
   email: ketua@institusi.ac.id\\
   \textsuperscript{2}Nama Fakultas, Nama Institusi, Kota, Kodepos, Negara,
   email: anggota1@institusi.ac.id\\
   \textsuperscript{3}Nama Fakultas, Nama Institusi, Kota, Kodepos, Negara,
   email: anggota2@institusi.ac.id\\
   \textsuperscript{4}Nama Fakultas, Nama Institusi, Kota, Kodepos, Negara\\[0.3em]
   \textit{Corresponding Author}: Nama Ketua Tim\par}
\end{center}
\vspace{1.0em}

{\small\justifying
% INTISARI: 200-300 kata, 4 bagian wajib (pengantar ~1/4 bagian, tujuan,
% metodologi singkat, penjelasan hasil), TANPA persamaan/tabel/sitasi.
\noindent\textbf{\textit{INTISARI}} --- [Pengantar/latar belakang singkat...]
[Tujuan penelitian...] [Metodologi singkat...] [Penjelasan hasil penelitian
secara keseluruhan, sertakan angka konkret jika ada.]
\par\vspace{0.7em}
% KATA KUNCI: 4-8 kata, dipisah koma, tiap kata diawali huruf kapital.
\noindent\textbf{KATA KUNCI:} Kata Kunci 1, Kata Kunci 2, Kata Kunci 3,
Kata Kunci 4, Kata Kunci 5.
\par}
\vspace{1.5em}
\end{@twocolumnfalse}
}]

%==============================================================================
% ISI (dua kolom)
%==============================================================================
\justifying

\section{Pendahuluan}
% Tulis sebagai narasi mengalir: latar belakang -> masalah spesifik ->
% kelemahan pendekatan sebelumnya (dengan sitasi) -> celah penelitian ->
% kontribusi -> (opsional) rumusan masalah -> manfaat/dampak strategis.
% Lihat references/common_juri_complaints.md poin 3 sebelum menulis ini.
[ISI PENDAHULUAN]

\section{Tinjauan Pustaka}
% Atau "Kajian Teori" kalau perlu subsection per konsep dasar.
% Pola tiap paragraf: "[Metode X] mencapai [metrik], namun [keterbatasan]
% [sitasi]." lalu simpulkan celah di akhir bagian.
[ISI TINJAUAN PUSTAKA]

\section{Metodologi}

\subsection{Dataset}
% WAJIB konkret: sumber, metode akuisisi, timeframe, frekuensi, jumlah
% baris/sampel, durasi pengumpulan, contoh satu baris data.
[ISI DATASET — gunakan tabel seperti contoh di bawah jika lebih dari satu
sumber data]

\begin{table}[t]
\centering
\caption{Jenis Data, Variabel, dan Sumber Dataset}
\label{tab:dataset}
\footnotesize
\begin{tabular}{@{}p{1.3cm}p{1.8cm}p{1.5cm}p{1.5cm}@{}}
\toprule
\textbf{Kategori} & \textbf{Variabel} & \textbf{Sumber} & \textbf{Frekuensi} \\
\midrule
{[Kategori]} & [Variabel] & [Sumber] & [Frekuensi] \\
\bottomrule
\end{tabular}
\end{table}

\subsection{Arsitektur/Metode Usulan}
[ISI METODE, rujuk ke Gambar~\ref{fig:arsitektur} untuk diagram alur]

\begin{figure}[t]
\centering
% Ganti \placeholderfig{...} dengan \includegraphics[width=0.92\linewidth]
% {nama_file.png} setelah pengguna menyediakan gambar aslinya.
\placeholderfig{7cm}{4cm}{Arsitektur sistem / diagram alur metode}
\caption{[Judul gambar arsitektur/alur sistem].}
\label{fig:arsitektur}
\end{figure}

\section{Hasil dan Diskusi}
% Jika eksperimen belum lengkap, nyatakan status secara eksplisit di sini
% (lihat references/common_juri_complaints.md poin 6). Jangan biarkan
% ambigu apakah hasil sudah final atau belum.
[ISI HASIL — setiap klaim angka harus merujuk ke tabel]

\begin{table}[t]
\centering
\caption{[Judul Tabel Hasil]}
\label{tab:hasil}
\footnotesize
\begin{tabular}{@{}lccc@{}}
\toprule
\textbf{Model/Skenario} & \textbf{Metrik 1} & \textbf{Metrik 2} & \textbf{Metrik 3} \\
\midrule
{[Baris 1]} & [Angka] & [Angka] & [Angka] \\
\bottomrule
\end{tabular}
\end{table}

\section{Kesimpulan}

\subsection{Kesimpulan}
[Rangkum hasil kuantitatif terkuat + jawab rumusan masalah/RQ jika ada]

\subsection{Saran}
[Rekomendasi konkret dan actionable untuk pekerjaan selanjutnya]

%==============================================================================
% REFERENSI (gaya IEEE, tanpa nomor heading, font 8pt)
%==============================================================================
\begin{thebibliography}{99}
\footnotesize
\bibitem{ref1}
[Penulis], ``[Judul],'' \textit{[Nama Jurnal/Prosiding]}, [Tahun].

\end{thebibliography}

\end{document}
